%! Logarithmic Function Context and Data Modeling — Unit 2, Lesson 2.14 — Group Activity (Answer Key)
\documentclass[10pt]{article}
\usepackage{precalculus-article}
\usepackage{precalculus-key}   % pulls in -boxes; do NOT also load -boxes
\newcommand{\TallMath}[1]{$\displaystyle #1\rule[-1.4em]{0pt}{3.2em}$}

\begin{document}

\pageheader{Unit 2, Lesson 2.14}{Group Activity --- Answer Key}
\namepartnerperiod

\vspace{0.1in}

\begin{scenariobox}[Species Diversity and Sound: Logarithmic Models in Action]{sky}
An ecologist studying habitat biodiversity and a physicist studying sound both encounter
data that are best modeled with logarithmic functions. Your group will construct, interpret,
and apply logarithmic models in each setting.
\end{scenariobox}

\vspace{0.1in}

\begin{tcolorbox}[colback=white, colframe=black!40,
    title=\textbf{Tier R --- Remediate},
    fonttitle=\bfseries, arc=2mm, left=3mm, right=3mm, top=2mm, bottom=2mm]

\textbf{Species Diversity --- reading a table and using a simple model.}

\vspace{4pt}
The table below shows how species count $S$ depends on habitat area $A$ (hectares).

\vspace{4pt}
\begin{center}
\renewcommand{\arraystretch}{1.8}
\begin{tabular}{|c|c|c|c|c|}
\hline
\rowcolor{sky}
$A$ & $1$ & $10$ & $100$ & $1000$ \\
\hline
$S$ & $0$ & $1$  & $2$   & $3$ \\
\hline
\end{tabular}
\end{center}

\begin{enumerate}[leftmargin=1.8em, itemsep=10pt]

\item As $S$ increases by 1, how does $A$ change?

\vspace{4pt}
$A$ \ans{multiplies by 10 (increases tenfold)}. The proportion factor is: \ans{10}

\item Write the logarithmic model.

\vspace{4pt}
$S(A) = $ \ans{$\log_{10}(A)$}

\item Use the model to predict $S(10000)$.

\vspace{4pt}
$S(10000) = $ \ans{$\log_{10}(10000) = 4$}

\item Is $S(0.5)$ defined? Explain.

\ansline{Yes, $S(0.5) = \log_{10}(0.5) \approx -0.301$. The domain of $\log_{10}(A)$ is
$A > 0$, and $0.5 > 0$, so the function value is defined (though negative, meaning area
less than 1 hectare gives a negative species count in this model, which may not be
contextually meaningful).}

\end{enumerate}
\end{tcolorbox}

\vspace{0.1in}

\begin{tcolorbox}[colback=white, colframe=black!40,
    title=\textbf{Tier A --- Approaching Proficiency},
    fonttitle=\bfseries, arc=2mm, left=3mm, right=3mm, top=2mm, bottom=2mm]

\textbf{Species Diversity --- building a model from two data points.}

\vspace{4pt}
Model form: $S(A) = a\log_3(A) + k$.

\begin{enumerate}[leftmargin=1.8em, itemsep=10pt]

\item System of equations:

\vspace{4pt}
Equation 1: \ans{$2 = a\log_3(1) + k$}

\vspace{4pt}
Equation 2: \ans{$4 = a\log_3(9) + k$}

\item Since $\log_3(1) = $ \ans{0}, Equation 1 gives $k = $ \ans{2}

\item Substituting into Equation 2:

\ansline{$4 = a\log_3(9) + 2 = 2a + 2 \;\Rightarrow\; 2 = 2a \;\Rightarrow\; a = 1$}

$a = $ \ans{1}

\item Complete model:

\vspace{4pt}
$S(A) = $ \ans{$\log_3(A) + 2$}

\item Predict $S(27)$:

\vspace{4pt}
$S(27) = $ \ans{$\log_3(27) + 2 = 3 + 2 = 5$}

\end{enumerate}
\end{tcolorbox}

\vspace{0.1in}

\begin{tcolorbox}[colback=white, colframe=black!40,
    title=\textbf{Tier E --- Extension},
    fonttitle=\bfseries, arc=2mm, left=3mm, right=3mm, top=2mm, bottom=2mm]

\textbf{Sound intensity --- transformations and real-world predictions.}

\vspace{4pt}
$S(I) = 10\log_{10}(I/I_0)$, $I_0 = 10^{-12}$ W/m$^2$.

\begin{enumerate}[leftmargin=1.8em, itemsep=10pt]

\item Expand:

\vspace{4pt}
$\log_{10}(10^{-12}) = $ \ans{$-12$}, so $-10\log_{10}(I_0) = $ \ans{$120$}.

\vspace{4pt}
$S(I) = 10\log_{10}(I) + $ \ans{120}

\item Vertical shift of $S(I) = 10\log_{10}(I) + 120$ versus $f(I) = 10\log_{10}(I)$:

\vspace{4pt}
Vertical shift: \ans{$+120$ (up 120)}

\item Find $I$ for $S = 60$:

\vspace{4pt}
$60 = 10\log_{10}(I) + 120$

\vspace{4pt}
$\log_{10}(I) = $ \ans{$-6$}

\vspace{4pt}
$I = $ \ans{$10^{-6}$} W/m$^2$

\item Find $I$ for $S = 120$, then the ratio:

\vspace{4pt}
$120 = 10\log_{10}(I) + 120 \;\Rightarrow\; \log_{10}(I) = 0 \;\Rightarrow\; I = $ \ans{$10^{0} = 1$} W/m$^2$

\vspace{4pt}
Ratio: \ans{$1 \div 10^{-6} = 10^6$} --- the rock concert is \ans{one million} times more intense.

\end{enumerate}
\end{tcolorbox}

\end{document}
